% Physics 310 -- Lab 3, themed alternative of "Phys310 - Lab 3 (2024)_redo.pdf".
%
% Same background, procedure and questions as the 2024 worksheet; only the
% presentation differs. Figures were extracted from that PDF into figures\.
% See ..\physics310.sty. REQUIRES LuaLaTeX:
%
%   powershell -ExecutionPolicy Bypass -File build-handout.ps1 -Tex labs\Lab3_Alt.tex -Engine lualatex

\documentclass[11pt]{article}
\usepackage[compact]{physics310}

% ---------------------------------------------------------------------------
% Lab metadata -- the only things that change from one lab to the next
% ---------------------------------------------------------------------------
\newcommand{\longname}{Lab 3}
\newcommand{\shortname}{Lab 3}
\newcommand{\term}{Fall 2026}
\newcommand{\totalpoints}{35}

\physicsheaders

% A gamma spectrum to identify, with a rule to write the isotope on.
\newcommand{\spectrum}[1]{%
    \par\addvspace{6pt}\needspace{2.9in}%
    \noindent\includegraphics[width=\linewidth]{figures/#1}\par
    \vspace{2pt}
    \noindent\hspace*{0.35in}\textbf{Isotope:}\ \fillin[2.2in]\par
    \addvspace{9pt}}

\begin{document}

\titleblock{Laboratory}
    {Radiation Detectors and Gamma Spectra \textperiodcentered\ United States Air Force Academy,
     Department of Physics and Meteorology}
    {\chip{\totalpoints\ Points}}

\labnames{3}

The purpose of this lab is to use detectors to determine various things about unknown radioactive
sources. You will first use a Geiger counter to locate a source, and then you will use a
scintillation detector to identify it.

\labsection{I}{Background}

\labsub{A}{Gamma-Ray Interactions with Matter}

The three dominant methods gamma-rays interact with material are the photoelectric effect, Compton
scattering, and electron-positron pair production.

In the photoelectric effect, the photon deposits 100\% of its energy into an electron, freeing it
from its bound state in an atom. When this occurs inside a detector, the output pulse produced will
be proportional to the initial energy since all of the energy is transferred to the material.

In Compton scattering, the photon can scatter off of an electron transferring only a portion of its
energy to the electron and retaining the rest. The amount of energy transferred depends on the
scattering angle where the maximum energy occurs at 180\textdegree. If the photon does not go on to
interact further with the detector, only a portion of the initial energy is measured by the
detector. However, this maximum energy can produce a noticeable feature in scintillation detectors,
known as the Compton edge, and can be calculated from the incident gamma ray energy. In pair
production, an electron-positron pair are created when a gamma-ray interacts with a nucleus of an
atom. These newly formed elementary particles slow down through scattering interactions, and, as
the positron comes to rest, it annihilates with an electron producing a pair of back-to-back
511 keV gamma-rays. All of these interactions have the possibility of being measured by the
detector. However, pair production can only occur if the photon has an energy greater than
1.022 MeV (the mass-energy of the electron-positron pair).

\labsub{B}{Scintillation Detectors}

You have already worked with Geiger counters in previous labs. In this lab, we will also be working
with a new type of detector --- a scintillation detector. Scintillation detectors are created using
materials that will scintillate, or produce light, when radiation deposits energy. They are one of
the oldest types of radiation detectors, since measurements could be made with photographic film or
the human eye. In modern scintillation detectors, light output is converted into current pulses
using a photomultiplier tube as shown in the figure. The most common type of modern scintillators
are made from thallium doped sodium iodide [NaI(Tl)] crystals which have ideal scintillation
properties. That is what we'll be using today.

\begin{center}
    \includegraphics[width=4.1in]{figures/lab3-p1-1.png}
\end{center}

One of the most useful features of modern scintillation detectors is their ability to produce a
light yield, and thus a current pulse, proportional to the energy deposited in the material by the
radiation at a high efficiency. This proportionality provides the user with valuable information
about the incident radiation's energy and therefore its source.

\labsub{C}{Gamma Spectra}

\noindent\begin{minipage}[c]{0.42\linewidth}
    \includegraphics[width=\linewidth]{figures/lab3-p2-1.png}
\end{minipage}\hfill
\begin{minipage}[c]{0.54\linewidth}
    \setlength{\parskip}{0pt}
    \begin{tabularx}{\linewidth}{@{}c >{\raggedright\arraybackslash}X l@{}}
        \rowcolor{teal}
        \thead{Feature} & \thead{Cs-137 Signature} & \thead{Energy} \\
        \rowcolor{card} \datarow a & Gamma Decay Peak & 662 keV \\
        \datarow b & Compton Edge & ?? \\
        \rowcolor{card} \datarow c & Backscatter & 184 keV \\
        \datarow d & X-ray Peak & 32 keV \\
    \end{tabularx}
\end{minipage}

\vspace{8pt}
In order to deduce meaningful information about a radiation source using a scintillation detector,
many individual measurements are required. If each current pulse is measured, binned, and plotted
as a histogram, features begin to emerge which can be very useful. For example, when we examine the
figure above --- if many current pulses are measured with the same magnitude, a large peak (a) is
produced in the histogram. This tells us that one particular energy has been deposited in the
scintillating material repeatedly, and likely indicates a radioactive source is present and
producing radioactive particles with that very specific energy. Because every radioactive isotope
has a unique decay mechanism, they each produce a unique energy signature. These signatures act as
an identifiable fingerprint for each isotope. By correlating the features revealed in the histogram
from the scintillation detector with known energy signatures, unknown isotopes can be identified.

\labsub{D}{Energy Spectrum Calibration}

The ability of scintillation detectors to produce these energy spectra are the primary
characteristic that makes them one of the most useful detectors available. However, the spectra
directly produced by the scintillation detectors are histograms of binned current pulse magnitudes
(also known as channels), and they must be converted to energy spectra before they can be used. In
order to convert the histogram from channels to energy, a calibration equation must be determined.

To accomplish this, a known source is measured with the detector first. We will use the Cs-137
source from the previous section as an example. Once a spectrum is gathered from the detector with
the known source, a correlation can be drawn from the corresponding channels of each peak to the
energies that are known to be produced by the known source. If channels for two features can be
gathered from the histogram and correlated to the same features of the energy signatures from the
isotope, the equation of a line can be found and used to convert any channel to its corresponding
energy for future isotopes.

\begin{equation}
    Energy = \frac{E_1 - E_2}{C_1 - C_2} * channel
\end{equation}

\labsub{E}{Equipment}

\begin{itemize}
    \item Pancake counter (Geiger-Mueller detector)
    \item FLIR Identifinder (scintillation detector)
    \item FLIR Spectrum: \url{http://422278800004.local/App/Spectrum}
\end{itemize}

\labsub{F}{Procedure and Data}

\begin{steps}
    \item Cs-137 is a common calibration source for field detectors. A sample Cs-137 spectrum will
          be shown at the beginning of class.

          \vspace{4pt}
          \begin{minipage}{\linewidth}
          \setlength{\parskip}{0pt}
          \noindent\begin{tabularx}{\linewidth}{@{}>{\raggedright\arraybackslash}p{1.8in} X X@{}}
              \rowcolor{teal}
              \thead{Known Feature} & \thead{Energy} & \thead{Channel Number} \\
              \rowcolor{card} \datarow --- & 0 keV & 0 \\
              \datarow X-ray Peak & 32 keV & \\
              \rowcolor{card} \datarow Compton Edge & & 160 \\
              \datarow Primary Decay Peak & 662 keV & 220 \\
          \end{tabularx}
          \end{minipage}

          \vspace{6pt}
          Use the channel number of the primary peak as given and the channel number of the Compton
          edge along with Eq (1) to estimate the channel number of the 32 keV X-ray peak and the
          energy of the Compton edge.
          \ptsbox{6}{1.9in}

    \item Get a pancake probe. Use it to survey the closed containers around the classroom and
          record the dose rate you measure at each. Do this until you find one that you believe
          contains a source (give a dose rate for at least 4 of the containers regardless of how
          quickly you find a source).

          Once you have picked one that you believe contains a source, alert your instructor who
          will open the box.

          \vspace{4pt}
          \llap{\chip[accent]{3 pts}\hspace{9pt}}%
          \begin{minipage}{\linewidth}
          \setlength{\parskip}{0pt}
          \noindent\begin{tabularx}{\linewidth}{@{}X X@{}}
              \rowcolor{teal}
              \thead{Container Description} & \thead{Measured Dose Rate} \\
              \rowcolor{card} \datarow & \\
              \datarow & \\
              \rowcolor{card} \datarow & \\
              \datarow & \\
              \rowcolor{card} \datarow & \\
              \datarow & \\
          \end{tabularx}
          \end{minipage}

    \item Analyze the spectra shown below for 5 different isotopes. Use your equation sheet to
          determine what spectra corresponds to what isotope.
          \ptsnote{10 pts}

          \vspace{2pt}
          \textit{Isotopes to consider: Na-24, K-40, Co-60, Mo-99, I-131, Xe-135, Cs-137, Am-241.}
\end{steps}

\spectrum{lab3-p4-1.png}
\spectrum{lab3-p4-2.png}
\spectrum{lab3-p4-3.png}
\spectrum{lab3-p5-1.png}
\spectrum{lab3-p5-2.png}

\labsection{G}{Discussion Questions}

\question[5] For the final spectrum, the two primary decay peaks were found at channels 390 and
    442. Use Eq (1) to numerically determine the energies of these peaks.
\writingbox{1.8in}

\question[5] The equation normally used to calculate the energy of the Compton edge is
\begin{equation*}
    E_{compton} = \frac{2E^2}{m_e c^2 + 2E}
\end{equation*}
where \(E\) is the energy of the initial gamma, and \(m_e c^2\) is the rest energy of the electron.
Calculate the Compton edge energy using the energy calculated above for the primary peak.
\writingbox{1.8in}

\question[3] Discuss the importance of scintillation detectors in radiation safety and nuclear
    science. What benefits do these detectors have over the Geiger-Mueller tubes?
\writingbox{1.6in}

\question[3] That said, your pancake counter (a gas-filled detector) was also useful in this lab.
    What are some advantages of gas counters? Identify a situation in which you might prefer one of
    these detectors over a scintillator or solid-state detector.
\writingbox{1.6in}

\end{document}
